\chapter{Точкови оценки}

\section{Постановка}

\begin{rembox}[Идея]
Нека $X$ описва някакво свойство на генерална съвкупност, с функция
на разпределение $F_X(x)=\mathbb P(X<x)$ (или плътност $f_X$).
Разполагаме с \textbf{извадка} $\bar X=(X_1,\dots,X_n)$, където
$(X_i)_{i=1}^n$ са i.i.d.\ (независими наблюдения над $X$,
$X_i\stackrel d=X$), и искаме на базата на извадката да определим
или приближим $F_X$ (респ.\ $f_X$). Важно е да отбележим: това, че
знаем $f_X$ или $F_X$, не означава, че познаваме случайния изход
$X$ предварително - знаем само разпределението му.
\end{rembox}

\begin{rembox}[Допускане: параметричен клас]
Допускаме, че $X$ принадлежи на някакъв \textbf{клас} от случайни
величини $F_X(x;\theta)$, където $\theta$ е параметър (или вектор
от параметри), асоцииран с възможните разпределения на $X$. Например:
\begin{itemize}[nosep]
  \item $\theta=p$: $X\sim Be(p)$, $p\in(0,1)$;
  \item $\theta=\lambda$: $X\sim Exp(\lambda)$, $F_X(x;\lambda)=1-e^{-\lambda x}$, $\lambda>0$;
  \item $\theta=(\mu,\sigma^2)$: $X\sim N(\mu,\sigma^2)$, $\mu\in\mathbb R$, $\sigma^2>0$.
\end{itemize}
\end{rembox}

\begin{defbox}[~(Точкова оценка)]
Нека $X$ е от класа $F_X(x;\theta)$ и нека $\bar X=(X_1,\dots,X_n)$
е извадка от независими наблюдения над $X$. Функцията
$\hat\theta=\hat\theta(\bar X)$ се нарича \textbf{точкова оценка}
(или \textbf{статистика}) за $\theta$. Тъй като $\hat\theta$ е
функция на случайни величини, самата $\hat\theta$ е случайна
величина; за конкретна реализация $\bar x=(x_1,\dots,x_n)$ получаваме
конкретно число $\hat\theta=\hat\theta(\bar x)$.
\end{defbox}

\section{Максимално правдоподобни оценки}

\begin{defbox}[~(Функция на правдоподобие)]
Нека $X$ е от класа $f_X(x;\theta)$ (или дискретна с $f(x;\theta)=
\mathbb P(X=x;\theta)$) и $\bar X=(X_1,\dots,X_n)$ е извадка от
независими наблюдения над $X$. \textbf{Функцията на правдоподобие}
е съвместната плътност на извадката:
\[
  L(\bar X;\theta) = \prod_{j=1}^{n} f_X(X_j;\theta).
\]
\end{defbox}

\begin{defbox}[~(Максимално правдоподобна оценка, МПО)]
Нека $\Theta$ е множеството от възможни стойности на $\theta$.
Оценката $\hat\theta$, за която
\[
  L(\bar X,\hat\theta) = \sup_{\theta\in\Theta} \prod_{j=1}^{n} f(X_j;\theta),
\]
се нарича \textbf{максимално правдоподобна оценка} за $\theta$. Т.е.\
$\hat\theta$ е стойността на параметъра, за която функцията на
правдоподобие достига максимум - интуитивно, тази стойност на
$\theta$, при която наблюдаваната извадка би била „най-вероятна“.
\end{defbox}

\begin{rembox}[Логаритмуване]
Ако $f_X$ е диференцируема по $\theta$, МПО обикновено се търси чрез
$\dfrac{\partial L}{\partial\theta}=0$ или (по-удобно на практика)
чрез $\dfrac{\partial\ln L}{\partial\theta}=0$ - понеже логаритъмът е
монотонна функция, $L$ и $\ln L$ достигат максимум в една и съща
точка.
\end{rembox}

\begin{rembox}[Пример: едно наблюдение от \texorpdfstring{$N(\mu,1)$}{N(мю,1)}]
Нека $X\sim N(\mu,1)$ и $\bar X=X_1=x_1$ (едно единствено
наблюдение). Тогава $L(\bar X,\mu)=f_X(x_1,\mu)$ - плътността на
$N(\mu,1)$, изчислена в точката $x_1$, като функция на $\mu$. Тя се
максимизира, когато центърът на камбанката съвпада точно с
наблюдението, т.е.\ $\hat\mu=x_1$ (интуитивно: измежду всички
възможни центрове $\mu$, стойността $\mu=x_1$ прави наблюдението
$x_1$ „най-вероятно“).
\end{rembox}

\begin{thmbox}[]{Теорема}
Нека $X\sim N(\mu,\sigma^2)$ и $\bar X=(X_1,\dots,X_n)$ са
независими наблюдения над $X$. Тогава МПО за $\mu$ е
\[
  \hat\mu = \overline{X}_n = \frac1n\sum_{j=1}^{n} X_j
\]
(винаги, независимо дали $\sigma^2$ е известно), а МПО за $\sigma^2$
зависи от това дали $\mu$ е известно:
\[
  \text{а) } \hat\sigma^2 = \frac1n\sum_{j=1}^n (X_j-\mu)^2\ \text{(ако $\mu$ е известно)},
\]
\[
  \text{б) } \hat\sigma^2 = \frac1n\sum_{j=1}^n (X_j-\overline X_n)^2\ \text{(ако $\mu$ не е известно)}.
\]
\end{thmbox}

\begin{proofbox}
Функцията на правдоподобие е
\[
  L(\bar X,\theta) = \Big(\frac1{\sqrt{2\pi}}\Big)^{n}(\sigma^2)^{-n/2}\,e^{-\sum_{j=1}^n\frac{(X_j-\mu)^2}{2\sigma^2}}, \qquad \theta=(\mu,\sigma^2),
\]
следователно
\[
  \ln L(\bar X,\theta) = -\frac n2\ln(2\pi) - \frac n2\ln(\sigma^2) - \sum_{j=1}^{n}\frac{(X_j-\mu)^2}{2\sigma^2}.
\]
Диференцираме по $\mu$ и приравняваме на нула:
\[
  \frac{\partial\ln L}{\partial\mu} = \frac1{\sigma^2}\sum_{j=1}^n (X_j-\mu) = 0 \quad\Longrightarrow\quad \hat\mu = \overline X_n
\]
(независимо от стойността на $\sigma^2$). Диференцираме по $\sigma^2$:
\[
  \frac{\partial\ln L}{\partial\sigma^2} = -\frac{n}{2\sigma^2} + \frac1{2\sigma^4}\sum_{j=1}^n(X_j-\mu)^2 = 0.
\]
Ако $\mu$ е известно, решаваме директно: $\hat\sigma^2 = \tfrac1n
\sum_{j=1}^n(X_j-\mu)^2$. Ако $\mu$ не е известно, заместваме
неговата МПО оценка $\hat\mu=\overline X_n$ (решавайки системата от
двете уравнения едновременно): $\hat\sigma^2=\tfrac1n\sum_{j=1}^n
(X_j-\overline X_n)^2$. \qedbg
\end{proofbox}

\begin{rembox}[Пример: \texorpdfstring{$X\sim U(0,\theta)$}{X~U(0,тета)}]
Нека $f_X(x)=\tfrac1\theta\,I_{[0,\theta]}(x)$, $\theta>0$.
Тогава
\[
  L(\bar X,\theta) = \prod_{j=1}^{n}\frac1\theta\,I_{[0,\theta]}(X_j) = \begin{cases}\dfrac1{\theta^n}, & \theta\ge\max_{1\le j\le n}X_j\\ 0, & \theta<\max_{1\le j\le n}X_j\end{cases}.
\]
Понеже $\tfrac1{\theta^n}$ намалява по $\theta$, функцията на
правдоподобие е максимална при най-малкото допустимо $\theta$, т.е.
\[
  \hat\theta = \max_{1\le j\le n} X_j.
\]
\end{rembox}

\section{Метод на моментите}

\begin{rembox}[Идея]
От усиления закон за големите числа (Тема~12), за i.i.d.\ извадка:
$\overline X_n^{(j)} := \tfrac1n\sum_{k=1}^n X_k^j
\xrightarrow{\text{п.с.}}\mathbb EX^j$ за всеки ред на момента $j$
(стига моментът да съществува). Идеята на метода на моментите е да
\emph{изравним} теоретичните моменти (функции на $\theta$) с
емпиричните моменти от извадката и да решим спрямо~$\theta$.
\end{rembox}

\begin{defbox}[~(Метод на моментите, ММ)]
Нека $X$ е случайна величина от клас $F_X(x;\theta)$, $\theta=
(\theta_1,\dots,\theta_s)$, и $\bar X$ е извадка. Ако първите $s$
момента на $X$ съществуват, полагаме $\mu^{(j)}(\theta):=\mathbb
EX^j$, $1\le j\le s$, и емпиричните моменти $\overline
X_n^{(j)}=\tfrac1n\sum_{k=1}^n X_k^j$. Тогава \textbf{оценката по
метода на моментите} $\hat\theta=(\hat\theta_1,\dots,\hat\theta_s)$
се получава чрез решаване на системата (спрямо целия вектор
$\theta$ едновременно - за $s\ge2$ всеки момент обикновено зависи от
\emph{няколко} от параметрите, така че координатно обръщане по
компоненти на $\theta$ най-често не е възможно)
\[
  \mu^{(j)}(\theta) = \overline X_n^{(j)}, \qquad 1\le j\le s.
\]
При $s=1$ системата е едно уравнение с едно неизвестно и решението
приема простия вид $\hat\theta = \big(\mu^{(1)}\big)^{-1}\big(\overline
X_n^{(1)}\big)$.
\end{defbox}

\begin{rembox}[Пример: \texorpdfstring{$X\sim U(0,\theta)$}{X~U(0,тета)}]
$\mu^{(1)}(\theta)=\mathbb EX=\tfrac\theta2$. Изравняваме с
$\overline X_n^{(1)}$: $\tfrac\theta2=\overline X_n^{(1)}
\Rightarrow \hat\theta = 2\overline X_n^{(1)} = 2\overline X_n$ -
различна от МПО оценката $\max_j X_j$, но също разумна оценка на
$\theta$.
\end{rembox}

\begin{rembox}[Пример: \texorpdfstring{$X\sim N(\mu,\sigma^2)$}{X~N(мю,сигма^2)}, двата параметъра неизвестни]
Тук $\theta=(\mu,\sigma^2)$, а системата е $\mathbb EX=\overline
X_n^{(1)}$, $\mathbb EX^2=\overline X_n^{(2)}$. От $\mathbb EX=\mu$
получаваме $\hat\mu=\overline X_n^{(1)}=\overline X_n$. От
$\mathbb EX^2 = DX+(\mathbb EX)^2=\sigma^2+\mu^2$ получаваме
\[
  \hat\sigma^2 = \overline X_n^{(2)} - \hat\mu^2 = \overline X_n^{(2)} - \big(\overline X_n^{(1)}\big)^2,
\]
което съвпада точно с МПО оценката от предната теорема (в случая,
когато $\mu$ не е известно) - в този пример методът на моментите и
методът на максималното правдоподобие дават \emph{една и съща}
оценка.
\end{rembox}

\section{Свойства на точковите оценки}

\begin{defbox}[~(Неизместеност)]
Казваме, че $\hat\theta$ е \textbf{неизместена} (несмесена) оценка
за $\theta$, ако $\mathbb E\hat\theta(\bar X) = \theta$ (респ.\
$\mathbb E\hat\theta_j = \theta_j$ за всяко $1\le j\le s$, ако
$\theta=(\theta_1,\dots,\theta_s)$). Величината $|\mathbb E\hat\theta
-\theta|$ се нарича \textbf{систематична грешка} на $\hat\theta$.
\end{defbox}

\begin{rembox}[Пример: \texorpdfstring{$\hat\mu=\overline X_n$}{мю-капа=X-чертата}, \texorpdfstring{$X\sim N(\mu,\sigma^2)$}{X~N(мю,сигма^2)}]
\[
  \mathbb E\hat\mu = \mathbb E\Big(\frac1n\sum_{j=1}^n X_j\Big) = \frac1n\sum_{j=1}^n \mathbb EX_j = \frac1n\cdot n\mu = \mu \quad\Longrightarrow\quad \hat\mu\ \text{е неизместена.}
\]
\end{rembox}

\begin{rembox}[Пример: дисперсията, ако \texorpdfstring{$\mu$}{мю} е известно]
$\hat\sigma^2=\tfrac1n\sum_{j=1}^n(X_j-\mu)^2$:
\[
  \mathbb E\hat\sigma^2 = \frac1n\sum_{j=1}^n\mathbb E(X_j-\mu)^2 = \frac1n\cdot n\sigma^2 = \sigma^2 \quad\Longrightarrow\quad \hat\sigma^2\ \text{е неизместена.}
\]
\end{rembox}

\begin{rembox}[Пример: дисперсията, ако \texorpdfstring{$\mu$}{мю} \textbf{не} е известно (важен случай!)]
Нека $\hat\sigma^2 = \tfrac1n\sum_{j=1}^n(X_j-\overline X_n)^2 =
\overline X_n^{(2)} - \big(\overline X_n^{(1)}\big)^2$ (разписано
чрез разкриване на квадрата, аналогично на $DX=\mathbb EX^2-
(\mathbb EX)^2$). Тогава
\[
  \mathbb E\hat\sigma^2 = \mathbb E\overline X_n^{(2)} - \mathbb E\big(\overline X_n^{(1)}\big)^2 = \mathbb EX_1^2 - \frac1{n^2}\,\mathbb E\Big(\sum_{j=1}^n X_j^2 + \sum_{i\neq j}X_iX_j\Big).
\]
Понеже $(X_i)$ са i.i.d., $\mathbb E\sum_j X_j^2 = n\,\mathbb EX_1^2$,
а за $i\neq j$ (независими, $\mathbb EX_iX_j=\mathbb EX_i\,\mathbb
EX_j=\mu^2$) има точно $n(n-1)$ наредени двойки, следователно:
\begin{align*}
  \mathbb E\hat\sigma^2 &= \mathbb EX_1^2 - \frac1{n^2}\Big(n\,\mathbb EX_1^2 + n(n-1)\mu^2\Big) = \Big(1-\frac1n\Big)\mathbb EX_1^2 - \frac{n-1}{n}\mu^2\\
  &= \frac{n-1}{n}\big(\mathbb EX_1^2-\mu^2\big) = \frac{n-1}{n}\,\sigma^2 \neq \sigma^2.
\end{align*}
Следователно тази оценка е \textbf{изместена}, със систематична
грешка $|\mathbb E\hat\sigma^2-\sigma^2|=\tfrac{\sigma^2}n$. Но
$\mathbb E\hat\sigma^2 = \tfrac{n-1}n\sigma^2 \to \sigma^2$ при $n\to\infty$ - казваме, че $\hat\sigma^2$
е \textbf{асимптотично неизместена}. \emph{(Затова на практика често
се предпочита нормираната с $\tfrac1{n-1}$ версия на извадковата
дисперсия - тя е точно неизместена.)}
\end{rembox}

\begin{rembox}[Пример: \texorpdfstring{$X\sim U(0,\theta)$}{X~U(0,тета)}, МПО оценката \texorpdfstring{$\hat\theta=\max_j X_j$}{тета-капа=max Xj}]
Намираме плътността на $\hat\theta$:
\[
  F_{\hat\theta}(x) = \mathbb P\Big(\max_{1\le j\le n}X_j\le x\Big) = \prod_{j=1}^{n}\mathbb P(X_j\le x) = \big(\mathbb P(X_1\le x)\big)^n = \begin{cases}0,&x\le0\\ (x/\theta)^n,&x\in(0,\theta)\\1,&x\ge\theta\end{cases}
\]
(използвахме независимостта на $(X_j)$), следователно $f_{\hat\theta}
(x) = \tfrac n{\theta^n}x^{n-1}$ за $x\in(0,\theta)$. Тогава
\[
  \mathbb E\hat\theta = \int_0^{\theta} x\cdot\frac n{\theta^n}x^{n-1}\,dx = \frac n{\theta^n}\int_0^\theta x^n\,dx = \frac n{\theta^n}\cdot\frac{\theta^{n+1}}{n+1} = \frac{n}{n+1}\,\theta < \theta,
\]
т.е.\ $\hat\theta$ е \textbf{изместена} оценка, но $\tfrac n{n+1}
\theta\to\theta$ при $n\to\infty$, т.е.\ тя е \textbf{асимптотично
неизместена}.
\end{rembox}

\begin{defbox}[~(Състоятелност)]
Нека $\hat\theta=\hat\theta(\bar X)$ е точкова оценка на $\theta$
(респ.\ $\theta=(\theta_1,\dots,\theta_s)$). Казваме, че $\hat\theta$
е \textbf{състоятелна} оценка на $\theta$, ако
\[
  \hat\theta \xrightarrow[n\to\infty]{\mathbb P} \theta
  \qquad
  \Big(\text{еквивалентно: }\forall\varepsilon>0:\ \mathbb P\big(|\hat\theta(\bar X)-\theta|>\varepsilon\big)\xrightarrow[n\to\infty]{}0\Big)
\]
(вероятността за грешка намалява, когато броят на наблюденията в
извадката расте).
\end{defbox}

\begin{defbox}[~(Силна състоятелност)]
Казваме, че $\hat\theta$ е \textbf{силно състоятелна} оценка на
$\theta$, ако $\hat\theta\xrightarrow[n\to\infty]{\text{п.с.}}\theta$.
От теоремата за видовете сходимост (Тема~12), силната състоятелност
влече състоятелност.
\end{defbox}

\begin{rembox}[Пример: \texorpdfstring{$\hat\mu=\overline X_n$}{мю-капа=X-чертата} е силно състоятелна]
По усиления закон за големите числа (Тема~12), директно приложен
към i.i.d.\ извадката $(X_j)$: $\overline X_n=\tfrac1n\sum_j X_j
\xrightarrow{\text{п.с.}}\mathbb EX_1=\mu$.
\end{rembox}

\begin{rembox}[Пример: \texorpdfstring{$\hat\sigma^2=\overline X_n^{(2)}-(\overline X_n^{(1)})^2$}{сигма-капа^2} е силно състоятелна]
По УЗГЧ, приложен поотделно към редиците $(X_j)$ и $(X_j^2)$:
$\overline X_n^{(1)}\xrightarrow{\text{п.с.}}\mathbb EX_1=\mu$ и
$\overline X_n^{(2)}\xrightarrow{\text{п.с.}}\mathbb EX_1^2$.
Функцията $g(a,b)=b-a^2$ е непрекъсната, следователно
\[
  \hat\sigma^2 = \overline X_n^{(2)}-\big(\overline X_n^{(1)}\big)^2 \xrightarrow[n\to\infty]{\text{п.с.}} \mathbb EX_1^2-\mu^2 = \sigma^2,
\]
т.е.\ $\hat\sigma^2$ е силно състоятелна (макар да е изместена за
крайно $n$, видяхме по-горе).
\end{rembox}
